\section{Limitações e Ameaças à Validade}
\label{sec:limitacoes}

O trabalho apresenta quatro limitações principais. Primeiro, a cobertura da coleta depende da infraestrutura digital das câmaras municipais, o que introduz viés regional e institucional. Segundo, a avaliação do tradutor ainda carece de conjunto de teste independente e análise mais sistemática dos casos em que a reescrita perde nuances importantes do texto legislativo original. O benchmark de recuperação já inclui um \textit{baseline} lexical BM25 e estimativas simples de incerteza por \textit{bootstrap}, mas ainda faltam \textit{baselines} híbridos, métodos de fusão e validação humana parcial dos erros mais críticos. Terceiro, embora agora exista um benchmark anotado de recuperação mais amplo, ele continua exploratório: cobre apenas 15 problemas, usa um pool derivado da união dos top-30 comparados, não estima \textit{recall} absoluto no corpus e foi julgado por um agente GPT-5.4, e não por especialistas humanos dos temas analisados. Isso reduz o custo de avaliação e fornece uma referência base padronizada para comparação entre métodos, mas não equivale a uma validação substantiva qualificada. Também faltam elementos importantes para uma seção de recuperação mais robusta, como mais problemas, candidatos adicionais fora do pool endógeno e validação estratificada com avaliadores humanos. Quarto, a análise com indicadores é observacional, sensível à qualidade temporal da base e insuficiente para sustentar inferência causal. A proxy educacional baseada em matrículas, o uso de homicídios no estudo de caso, o agrupamento textual por Jaccard e o ruído residual da base devem ser entendidos como simplificações de uma avaliação inicial. Em uso institucional, o sistema deve funcionar apenas como apoio exploratório, com validação humana e contextualização local antes de qualquer adoção, adaptação ou priorização administrativa.
